\chapter{P4：把评论变成可以反复查询的证据工具}

\section{承上启下}

P3 已经说清楚什么算证据，但每次都从大量原始文件翻找，运营人员和后续 Agent 都会很慢，也容易引用错商品。P4 要做一个本地“证据账本”：输入同样的问题，尽量得到同样的可回查证据。

\section{业务问题、输入与边界}

业务问题是：怎样快速找到某个商品、某个体验维度的支持或反对信号？输入是 P1 冻结数据、P2 选定范围、P3 证据合同和本地原始文件。输出是受控的本地证据服务，不是自动生成商业结论。

业务价值是减少重复翻评论的时间，并让运营人员知道结论来自哪里；技术价值是学习表关系、全文检索、稳定引用和失败时不发布半成品。Agent 不能直接读取任意原始文件，也不能在这里拥有“自己解释证据”的权限。

\section{开发过程}

\subsection{1. 开工前先核对上游}

构建前再次核对 P1--P3 的数据身份、范围和合同。可以把这一步理解为开工前核对“账本版本”。有一项不一致就停止，不能用目录里看起来相似的文件补救。

\subsection{2. 建立本地证据账本}

用商品表保存允许查询的商品范围，用评论表保存结构化文本和来源位置。SQLite 是一个可以放在本机上的小型数据库，先把它理解为“项目自己的账本”。本地 evidence ID 是定位键，不是官方 Amazon review id。

\subsection{3. 给账本加快速查词的索引}

FTS5 是 SQLite 的全文检索功能。它像给仓库货架贴上关键词索引，让程序能按商品和 aspect 找到相关文本。工具参数由本地代码检查，Agent 不能提交任意 SQL，也不能把一个商品的证据混到另一个商品。

\subsection{4. 先在临时账本里验收，再替换旧账本}

程序先写临时 SQLite 文件，关闭数据库连接和文件句柄，检查商品数、关联评论数、FTS 行数和内容摘要都正确后，才替换旧库。这叫\textbf{原子发布}：要么完整换成新版本，要么继续使用旧版本，不留下半成品。Windows 上曾经因为句柄未关闭导致替换失败，所以“显式关闭”是必要动作，不是多余形式。

\section{本阶段学到的知识}

\begin{itemize}
  \item SQLite：本地的结构化证据账本。
  \item FTS5：账本上的全文索引，帮助快速找到关键词。
  \item 精确 join：按明确键连接，避免商品串线。
  \item 稳定主键：每次回查都有固定门牌号。
  \item 原子发布：完整版本才替换，避免半成品被使用。
\end{itemize}

\begin{beginnerbox}
想象两个商品都出现“noise”。如果 evidence ID 只由文本生成，两个商品可能得到难以区分的定位；把数据版本和来源位置也绑定进去，系统才能知道它属于哪一批、哪一个商品、哪一条来源记录。
\end{beginnerbox}

\section{问题、修复与验证}

问题一是构建中途失败却留下半成品；用临时库、关闭句柄和原子替换修复。问题二是内部定位键出现在公开回执；用聚合结果和扫描门修复。验证包括上游身份、精确商品关联、FTS 行、引用回查和公开内容扫描。本地稳定 ID 不被表述为平台官方标识。

\section{阶段产出与验收边界}

本地证据库被 Git 忽略，公开只保留\filelink{artifacts/store-summaries/phase-4-build.json}和\filelink{docs/reports/phase-4-local-store.md}等聚合回执。P4 在其范围内为\statuspassed：冻结范围内的本地构建和回查门通过；不说明上游评论绝对正确。

\section{怎样进入下一阶段}

P5 会用固定的运营案例比较简单词法检索和 BM25 等方案，判断增加复杂度是否真的有帮助。P4 解决“能不能稳定取证”，还没有解决“哪种检索对运营任务更有用”。

\section{面试怎么讲}

\begin{interviewbox}
“我没有让模型直接读原始评论，而是把 P3 规则做成本地 SQLite/FTS5 证据工具。构建前检查上游身份，按明确键关联商品，在临时库中验证后原子替换，引用最后再回账本复核。这样解决的是可追踪和可复现，不把本地 ID 冒充官方 review ID。”
\end{interviewbox}

\textbf{掌握标准：}能画出数据版本、合同、商品表、评论表、FTS5、查询工具和引用校验的关系，并解释为什么失败构建不能直接覆盖旧库。
